% Original: PDF strana 8, gornji slajd; drugi deo odobrene podele.
\slajdid{02-s015b}
\begin{frame}[label=02-s015b]{Povezivanje uslova u proizvod zbirova}
\begin{columns}[T,onlytextwidth]
\begin{column}{.66\textwidth}
% element: 02-s015b:e01
\Oznaka{Sledeći red: \((C,B,A)=(0,0,1)\)}\vspace{2mm}
Izlaz je 0 ako je \(C=0\) \alert{I} \(B=0\) \alert{I} \(A=1\). Odgovarajući činilac je
\[\textcolor{DEcrvena}{\overline{\bar C\bar BA}=C+B+\bar A}.\]
\par\vspace{2mm}
% element: 02-s015b:e02
Na isti način dobijamo činilac za svaki od pet redova sa \(F=0\).
\par\vspace{3mm}
\Rezultat{Sve činioce povezujemo operacijom \alert{I}: izlaz I kola je 0 ako je \alert{bar jedan} njegov ulaz 0.}
\end{column}
\begin{column}{.29\textwidth}\centering
\Oznaka{Funkcionalna tabela}\vspace{3mm}
{\fontsize{9}{11}\selectfont\begingroup
\setlength{\tabcolsep}{2pt}
\renewcommand{\arraystretch}{1.04}
\par\noindent
\begin{tabularx}{\linewidth}{*{4}{>{\centering\arraybackslash}X}}
\toprule C&B&A&F\\\midrule
\textcolor{DEcrvena}{0}&\textcolor{DEcrvena}{0}&\textcolor{DEcrvena}{0}&\textcolor{DEcrvena}{0}\\
\textcolor{DEcrvena}{0}&\textcolor{DEcrvena}{0}&\textcolor{DEcrvena}{1}&\textcolor{DEcrvena}{0}\\
\textcolor{DEcrvena}{0}&\textcolor{DEcrvena}{1}&\textcolor{DEcrvena}{0}&\textcolor{DEcrvena}{0}\\
0&1&1&1\\
\textcolor{DEcrvena}{1}&\textcolor{DEcrvena}{0}&\textcolor{DEcrvena}{0}&\textcolor{DEcrvena}{0}\\
1&0&1&1\\
1&1&0&1\\
\textcolor{DEcrvena}{1}&\textcolor{DEcrvena}{1}&\textcolor{DEcrvena}{1}&\textcolor{DEcrvena}{0}\\\bottomrule
\end{tabularx}\par
\endgroup
}
\end{column}
\end{columns}
\par\vspace{2mm}
% element: 02-s015b:e03
\[\textcolor{DEcrvena}{F=(C+B+A)(C+B+\bar A)(C+\bar B+A)(\bar C+B+A)(\bar C+\bar B+\bar A)}\]
\end{frame}
